\scanpage{23}
$A$ anneau local de dim $n$, $\mathrm I$ un module dualisant
$\Omega$ un module fondamental pour $A$, avec un isom donné
\[
\boxed{\unsure{g|\eta|\mathfrak D}_1:H^n(A)\xrightarrow{\sim}D(\Omega)}
\]
correspondant à un accouplement
\[
\begin{aligned}[b]&\text{[où \deleted{\unsure{un}}}\end{aligned}\quad
\boxed{\begin{aligned}[b]
&H^n(A)\otimes\Omega=H^n(\Omega)\xrightarrow{g}\mathrm I\\
&g\in D(H^n(\Omega))
\end{aligned}}\quad\begin{aligned}[b]&\text{\deleted{\unsure{est}} ]}\end{aligned}
\]
et à un isom
\[
g_2:\widehat\Omega\xrightarrow{\sim}D(H^n(A))
\]

On a pour \uline{tout $M$ de type} fini
\[
\boxed{H^n(M)\xleftarrow[\sim]{\mathrm{can}}H^n(A)\otimes M}
\xrightarrow[\sim]{g_1\otimes\operatorname{id}_M}D(\Omega)\otimes M
\xrightarrow[\sim]{\mathrm{can}}D(\operatorname{Hom}(M,\Omega))
\]
d'où
\[
\begin{aligned}[b]&\text{et par là}\end{aligned}\quad
\boxed{\begin{aligned}[b]
g_{1,M}:{}&H^n(M)\xrightarrow{\sim}D(\operatorname{Hom}(M,\Omega))\\
g_{2,M}:{}&\operatorname{Hom}(M,\Omega)^{\wedge}\xrightarrow{\sim}D(H^n(M))
\end{aligned}}
\]
qui pour $M\simeq A$ se réduisent à $g_1,g_2$. Pour $M=\Omega$, on trouve
\[
\begin{aligned}
g_{1,\Omega}:{}&H^n(\Omega)\xrightarrow{\sim}D(\operatorname{Hom}(\Omega,\Omega))\\
g_{2,\Omega}:{}&\operatorname{Hom}(\Omega,\Omega)^{\wedge}\xrightarrow{\sim}D(H^n(\Omega))
\end{aligned}
\]
Enfin, on constate que l'élément identique de
$\operatorname{Hom}(\Omega,\Omega)^{\wedge}$ correspond à l'élément $g$ de $D(H^n(\Omega))$.
Par suite, le transposé de $g$ est l'hom
\[
{}^tg:\widehat A\longrightarrow\operatorname{Hom}(\Omega,\Omega)^{\wedge}
\]
qui transforme $1_{\widehat A}$ en $\operatorname{id}_{\widehat\Omega}$, donc $\lambda$ en l'homothétie $\lambda$.

\uline{Proposition} 9 L'annulateur de $\Omega$ est l'intersection $\unsure{\mathfrak N|N}$
des idéaux primaires associés aux idéaux pr. de
corang $n$ dans $A$, et $\Omega$ est aussi un
module fondamental pour $A'=A/\mathfrak N$. Les idéaux
pr. associés à $\Omega$ sont les idéaux pr. de corang $n$ de $A$. \unsure{$p$|$\mathfrak p$}
\margin{La longueur $\unsure{\mathrm{lg}_p|\mathrm l_p|\ell_{\mathfrak p}}(\Omega)=\mathrm{lg}_p(A)$ pour un tel $p$}

L'annulateur de $\Omega$ est celui de $D(\Omega)=H^n(A)$,
donc \deleted{\unsure{notre|toute}} la première assertion suit du th. 3; la
% archive sheet number: 22

